\documentclass[11pt]{article}
\usepackage[T1]{fontenc}
\usepackage{lmodern}
\usepackage{amsmath,amssymb,amsthm,mathtools}
\usepackage[margin=1in]{geometry}
\usepackage[hidelinks]{hyperref}

\newtheorem{theorem}{Theorem}
\newtheorem{proposition}{Proposition}
\newtheorem{remark}{Remark}
\DeclareMathOperator*{\argmin}{arg\,min}
\newcommand{\TT}{\mathbb{T}}
\newcommand{\ZZ}{\mathbb{Z}}
\newcommand{\calA}{\mathcal{A}^{\star}}
\newcommand{\calH}{\mathcal{H}_{\calA}}
\newcommand{\PA}{P_{k_{\star}}}
\newcommand{\ii}{\mathrm{i}}

\title{E59: Idempotence and Orthogonality of the Dominant-Annulus Fourier Projection}
\author{Richard J. Reyes}
\date{October 26, 2026}

\begin{document}
\maketitle

\begin{abstract}
We show that restriction to a fixed dominant Fourier annulus on the periodic domain
is an idempotent, self-adjoint contraction on $L^{2}$, hence an orthogonal projection.
We also establish its exact energy decomposition, its best-approximation property,
and the corresponding discrete Fourier transform identity.  The result concerns a
fixed sharp spectral mask; it does not assert invariance of the selected subspace
under nonlinear evolution.
\end{abstract}

\section{Definitions and theorem}
Let $d\geq 1$, let $\TT^{d}=(\mathbb{R}/2\pi\ZZ)^{d}$, and equip
$L^{2}(\TT^{d};\mathbb{C})$ with the normalized inner product
\begin{equation}
\langle A,B\rangle := \frac{1}{(2\pi)^{d}}
  \int_{\TT^{d}} A(x)\overline{B(x)}\,dx,
\qquad \|A\|_{2}^{2}=\langle A,A\rangle.
\end{equation}
Its Fourier coefficients are
\begin{equation}
\widehat A_{k}=\frac{1}{(2\pi)^{d}}
  \int_{\TT^{d}} A(x)e^{-\ii k\cdot x}\,dx,
\qquad k\in\ZZ^{d},
\end{equation}
and $A=\sum_{k\in\ZZ^{d}}\widehat A_{k}e^{\ii k\cdot x}$
with convergence in $L^{2}$.  Parseval's identity states
\begin{equation}\label{eq:parseval}
\langle A,B\rangle = \sum_{k\in\ZZ^{d}}
    \widehat A_{k}\overline{\widehat B_{k}},
\qquad \|A\|_{2}^{2}=\sum_{k\in\ZZ^{d}}|\widehat A_{k}|^{2}.
\end{equation}
Fix a preferred wavenumber $k_{\star}\geq0$ and bandwidth
$\Delta k\geq0$.  Define the \emph{fixed} dominant annulus
\begin{equation}\label{eq:annulus}
\calA:=\left\{k\in\ZZ^{d}:\bigl||k|-k_{\star}\bigr|\leq\Delta k\right\}.
\end{equation}
Let $\chi_{\calA}$ be its indicator, and set
\begin{equation}\label{eq:projection}
(\PA A)(x)
 :=\sum_{k\in\calA}\widehat A_{k}e^{\ii k\cdot x}
 =\sum_{k\in\ZZ^{d}}\chi_{\calA}(k)\widehat A_{k}e^{\ii k\cdot x}.
\end{equation}
The corresponding band-limited subspace is
\begin{equation}
\calH:=\left\{B\in L^{2}(\TT^{d}):
\widehat B_{k}=0\ \text{for all }k\notin\calA\right\}.
\end{equation}

\begin{theorem}[E59: dominant-annulus projection]\label{thm:e59}
The operator $\PA:L^{2}(\TT^{d})\to L^{2}(\TT^{d})$
defined by \eqref{eq:projection} is bounded and satisfies
\begin{equation}\label{eq:main}
\boxed{\PA^{2}=\PA,\qquad \PA^{*}=\PA,\qquad
\|\PA A\|_{2}\leq\|A\|_{2}\quad(A\in L^{2}).}
\end{equation}
It is the orthogonal projection onto $\calH$, with
\begin{equation}\label{eq:pythagoras}
\boxed{\|A\|_{2}^{2}=\|\PA A\|_{2}^{2}
              +\|(I-\PA)A\|_{2}^{2}.}
\end{equation}
Furthermore, $\PA A$ is the unique best $L^{2}$ approximation
to $A$ from $\calH$:
\begin{equation}\label{eq:best}
\boxed{\PA A=\argmin_{B\in\calH}\|A-B\|_{2}.}
\end{equation}
If $\calA\neq\varnothing$, then $\|\PA\|_{\mathrm{op}}=1$;
if $\calA=\varnothing$, then $\PA=0$.
\end{theorem}

\begin{proof}
\textbf{Well-definedness and boundedness.}
Because $\chi_{\calA}(k)\in\{0,1\}$, the coefficient sequence
$(\chi_{\calA}(k)\widehat A_k)_{k\in\ZZ^d}$ belongs to $\ell^{2}$.
Thus \eqref{eq:projection} defines an $L^{2}$ function.  By Parseval,
\begin{equation}\label{eq:contraction}
\|\PA A\|_{2}^{2}
 =\sum_{k\in\ZZ^{d}}\chi_{\calA}(k)^{2}|\widehat A_{k}|^{2}
 =\sum_{k\in\calA}|\widehat A_{k}|^{2}
 \leq\sum_{k\in\ZZ^{d}}|\widehat A_{k}|^{2}
 =\|A\|_{2}^{2}.
\end{equation}
In particular $\PA$ is a bounded linear map.

\medskip\noindent
\textbf{Idempotence.}
Uniqueness of Fourier coefficients gives
\begin{equation}
\widehat{(\PA A)}_k=\chi_{\calA}(k)\widehat A_k.
\end{equation}
Applying $\PA$ once more,
\begin{align}
\PA^{2}A
 &=\sum_{k\in\ZZ^{d}}\chi_{\calA}(k)
       \widehat{(\PA A)}_k e^{\ii k\cdot x}\\
 &=\sum_{k\in\ZZ^{d}}\chi_{\calA}(k)^{2}
       \widehat A_k e^{\ii k\cdot x}\\
 &=\sum_{k\in\ZZ^{d}}\chi_{\calA}(k)
       \widehat A_k e^{\ii k\cdot x}
 =\PA A.
\end{align}
The penultimate equality uses precisely
$\chi_{\calA}(k)^{2}=\chi_{\calA}(k)$ for every $k$.
Since $A$ is arbitrary, $\boxed{\PA^{2}=\PA}$.

\medskip\noindent
\textbf{Self-adjointness and range.}
For $A,B\in L^{2}$, Parseval yields
\begin{align}
\langle\PA A,B\rangle
 &=\sum_k\chi_{\calA}(k)\widehat A_k\overline{\widehat B_k}\\
 &=\sum_k\widehat A_k\overline{\chi_{\calA}(k)\widehat B_k}
 =\langle A,\PA B\rangle,
\end{align}
because $\chi_{\calA}$ is real.  Hence $\PA^{*}=\PA$.
Every $\PA A$ has Fourier support in $\calA$, so
$\operatorname{Ran}(\PA)\subseteq\calH$.
Conversely, if $B\in\calH$, then $\PA B=B$ by the multiplier
formula, so $\calH\subseteq\operatorname{Ran}(\PA)$.
Thus $\operatorname{Ran}(\PA)=\calH$.  Self-adjointness and
idempotence show that $\PA$ is its orthogonal projection.

\medskip\noindent
\textbf{Energy decomposition.}
The Fourier supports of $\PA A$ and $(I-\PA)A$ are disjoint,
so their inner product vanishes by Parseval:
$\langle\PA A,(I-\PA)A\rangle=0$.
Expanding $A=\PA A+(I-\PA)A$ gives \eqref{eq:pythagoras}.

\medskip\noindent
\textbf{Best approximation.}
For any $B\in\calH$, the vectors $(I-\PA)A$ and
$\PA A-B\in\calH$ are orthogonal.  Therefore
\begin{equation}
\|A-B\|_{2}^{2}
 =\|(I-\PA)A\|_{2}^{2}+\|\PA A-B\|_{2}^{2}
 \geq\|(I-\PA)A\|_{2}^{2},
\end{equation}
with equality exactly when $B=\PA A$.  This proves \eqref{eq:best}.
Finally, \eqref{eq:contraction} implies $\|\PA\|_{\mathrm{op}}\leq1$.
If $k\in\calA$, the nonzero function $e^{\ii k\cdot x}$ is fixed
by $\PA$, proving $\|\PA\|_{\mathrm{op}}=1$; when $\calA$ is empty,
\eqref{eq:projection} gives the zero operator.
\end{proof}

\section{Spectral concentration (E62)}
For $A\neq0$, define the retained spectral-energy fraction
\begin{equation}
\eta(A):=\frac{\sum_{k\in\calA}|\widehat A_k|^{2}}
                 {\sum_{k\in\ZZ^d}|\widehat A_k|^{2}}
       =\frac{\|\PA A\|_{2}^{2}}{\|A\|_{2}^{2}}.
\end{equation}
Equation \eqref{eq:contraction} immediately implies
\begin{equation}
0\leq\eta(A)\leq1,
\qquad
\eta(A)=1\quad\Longleftrightarrow\quad \PA A=A.
\end{equation}

\section{Discrete FFT implementation}
Let $F\in\mathbb{C}^{N\times N}$ be a unitary discrete Fourier
transform matrix ($F^{*}F=FF^{*}=I$), and let
$D=\operatorname{diag}(m_1,\ldots,m_N)$ with $m_j\in\{0,1\}$
represent a \emph{fixed} discrete annulus mask.  The physical-space
operator $P=F^{*}DF$ satisfies
\begin{equation}
P^{2}=F^{*}DF F^{*}DF=F^{*}D^{2}F=F^{*}DF=P,
\end{equation}
since $D^{2}=D$.  Likewise $P^{*}=P$.  Floating-point FFT
implementations satisfy these identities up to rounding error when
the same mask and transform conventions are used twice.

\section{Scope and qualifications}
\begin{remark}[Correct annulus convention]
The symmetric condition $\bigl||k|-k_{\star}\bigr|\leq\Delta k$
defines a shell about the preferred wavenumber.  The one-sided condition
$|k|-k_{\star}\leq\Delta k$ defines a different set containing low
wavenumbers.  Idempotence holds for a fixed mask of either kind, but
only the symmetric condition encodes the intended annulus.
\end{remark}

\begin{remark}[Smooth multipliers]
For any bounded Fourier multiplier $Q$ defined by
$\widehat{QA}_k=w(k)\widehat A_k$, one has $Q^{2}=Q$ if and only if
$w(k)^{2}=w(k)$ for every admissible Fourier mode $k$, i.e.,
$w(k)\in\{0,1\}$.  A genuinely tapered band-pass filter is therefore
not generally idempotent.
\end{remark}

\begin{remark}[Adaptive selections and nonlinear evolution]
The proof assumes that $k_{\star}$ and $\Delta k$ are held fixed
under successive applications.  If a ``dominant annulus'' is selected
afresh from each input, the resulting map need not be idempotent.
Moreover, the fact that $\PA$ is a projection does not imply that
$\calH$ is invariant under an unprojected nonlinear PDE, nor does it
by itself demonstrate localization, confinement, or stability.
\end{remark}

\section*{Conclusion}
For the fixed, sharp dominant-annulus Fourier mask used in E59,
\[
\boxed{\PA^{2}=\PA,\qquad\PA^{*}=\PA,\qquad
       \|\PA A\|_{2}\leq\|A\|_{2}.}
\]
Thus E59 is proved as an exact Hilbert-space operator statement.

\end{document}